% Results: the layout experiments. Why they look this way is in analysis.tex.
All latencies are medians over 15 runs, from DRAM (4M rows $\times$ 16 attributes, 512\,MB per
layout) unless the table size is the variable.

\subsection{Projectivity}
\label{sec:res:proj}

\begin{figure}[t]
\centering
\fig{projectivity}
\caption{Latency vs the number of attributes read, without a predicate. Left: SUM of $k$
attributes. Right: projection of $k$ attributes into row-major output tuples.}
\label{fig:projectivity}
\end{figure}

\Cref{fig:projectivity} varies the number $k$ of the 16 attributes a query reads.
\begin{itemize}
\item \textbf{Aggregation.} DSM's time grows in proportion to $k$, from 1.0\ms{} at $k{=}1$ to
18\ms{} at $k{=}16$, because it reads exactly $k$ arrays. NSM's hardly depends on $k$
(13--22\ms), because every row's cache lines come in anyway. At $k{=}1$ DSM is 15$\times$
faster, close to the 16$\times$ ratio of bytes moved. PAX reads only the needed minipages too,
but pays a fixed cost per minipage: 3.9\ms{} at $k{=}1$, 4$\times$ DSM (\cref{sec:res:pax}).
\item \textbf{Projection.} Every layout also writes $8k$ bytes of output per row. DSM still wins
at low projectivity (4.0\ms{} vs NSM's 19\ms{} at $k{=}1$) and up to $k{=}8$, and NSM wins at
$k{=}16$ (61\ms{} vs DSM's 72\ms{} and PAX's 79\ms). The crossover lies between $k{=}8$ and
$k{=}16$.
\item \textbf{Strategy.} Without a predicate, late is slower than early on DSM for $k \ge 2$,
by 30\% at $k{=}16$ (94 vs 72\ms), because it writes each vector's output one attribute at a
time (\cref{sec:analysis:recon}).
\end{itemize}

\subsection{Selectivity}
\label{sec:res:sel}

\begin{figure}[p]
\centering
\fig{selectivity}
\caption{Latency vs selectivity for every layout and strategy. Each row of panels is one
query; the predicate is \code{c0 < x} and $c_0$ is among the attributes read. Both axes are
logarithmic.}
\label{fig:selectivity}
\end{figure}

\begin{table}[t]
\centering
\caption{Selectivity at which late stops beating early, log-interpolated between measured
points.}
\label{tab:crossovers}
\gen{selectivity_crossovers}
\end{table}

\Cref{fig:selectivity} sweeps the selectivity of \code{c0 < x} from 0.01\% to 100\% for four
queries under the four fixed strategies, and \cref{tab:crossovers} gives where late stops
beating early.
\begin{itemize}
\item \textbf{DSM is where the strategy matters most.} At 0.01\%, late reads $c_0$ and a few
lines of the other attributes: \code{SELECT *} takes 2.3\ms{} against early's 21\ms, 9$\times$
less. Early's time is nearly flat up to about 5\%, because it always reads all $k$ attributes.
Late's grows with the selectivity and crosses early's at 2.4\% for $k{=}4$ and 7.4\% for
$k{=}16$: the more attributes late can skip, the longer it stays ahead. With $k{=}2$ there is
little to skip, and the two are within 2\% of each other from 0.5\% to 2\%. At 100\%, late is
45\% slower (99 vs 68\ms).
\item \textbf{NSM barely cares.} Every row's lines are read whatever the strategy, so the
strategies differ only in CPU work (15--23\ms{} for \code{SELECT *} at low selectivity).
\item \textbf{PAX early is slow at every selectivity} with the default 16\,KB pages:
\code{SELECT *} takes 55--63\ms{} up to 20\%, 2.6$\times$ DSM early, while PAX late follows DSM
late at a fixed offset (\cref{sec:analysis:cache}).
\item \textbf{Branches.} \emph{early-pipe} and \emph{late-branch} peak near 50\%, where the
branch is least predictable: for $k{=}2$ on DSM, late-branch takes 17.6\ms{} against late's
6.9\ms. The branch-free variants have no peak, so we use them from here on.
\item \textbf{Aggregation} has lower crossovers (DSM 1.4\%, PAX 3.3\%): early SUM writes
nothing and stays at about 3.6\ms{} on DSM, so late has less to save.
\item \textbf{Other predicates} behave like their selectivity. An equality is a range of width
one and costs what \code{c0 < x} costs at the same selectivity (7.5 vs 6.9\ms{} for DSM late
\code{SELECT *} at 1\%). A conjunction costs what its combined selectivity costs, because late
refines the selection vector with the second predicate only at the surviving positions: adding
\code{c1 = 7} to \code{c0 = 42} cuts DSM late from 7.5 to 3.0\ms.
\end{itemize}

\subsection{Table width}
\label{sec:res:width}

Keeping the table at 512\,MB while $N$ grows from 2 to 256 (so rows get wider and fewer)
changes the picture in two ways. For SUM(c0), DSM's time falls as $1/N$, while NSM's falls much
more slowly (25 to 2.1\ms), because beyond $N{=}8$ it still reads one line per row. For
\code{SELECT *} of every row, reconstruction gets more expensive for the columnar layouts as
rows widen (DSM 63\ms{} at $N{=}4$, 106\ms{} at $N{=}256$; PAX 68 to 131\ms), while NSM's
row-major copy stays at 66--69\ms. Beyond about 16 streams, neither the 8-way L1 nor the
prefetcher keeps up with stitching $N$ columns. At 1\% selectivity NSM becomes the fastest from
$N{=}128$ (4.8 vs 5.2\ms{} for DSM): each qualifying row is then a short contiguous range,
while DSM fetches one line per qualifying row from each of $N$ columns.

\subsection{Working-set size}
\label{sec:res:ws}

\begin{figure}[t]
\centering
\fig{workingset}
\caption{Time per row vs table size ($N{=}16$), from 16\,KB (in L1) to 512\,MB. Dotted lines
mark the cache capacities. Small tables are scanned repeatedly until a sample lasts 20\ms.}
\label{fig:workingset}
\end{figure}

\Cref{fig:workingset} varies the table size from 16\,KB to 512\,MB.
\begin{itemize}
\item \textbf{SUM(c0)} costs DSM 0.08--0.10\,ns per row from L2 and L3 and 0.26\,ns from
DRAM, and NSM 0.5--0.6\,ns in cache and 3.8\,ns from DRAM. The NSM/DSM ratio grows from about
5$\times$ in L2 to 15$\times$ in DRAM: in cache the difference is CPU work (a strided load per
row against vectorized contiguous loads), in DRAM it is bytes.
\item \textbf{SUM of all 16 attributes.} From DRAM every layout moves the same 128 bytes per
row (4.3--5.5\,ns). In cache NSM is twice as slow as DSM (2.5 vs 1.1--1.4\,ns), because summing
across a row vectorizes worse than summing down a column.
\item \textbf{\code{SELECT *}.} DSM is the fastest while the table fits in L2 or a quarter of
L3 (2.2--3.1\,ns per row against NSM's 3.1--3.5), and NSM once the table fills L3 (13.0 vs
13.6\,ns at 16\,MB, 16.1 vs 18.1\,ns from DRAM). The reconstruction crossover depends on
where the data lives, not only on the layout.
\item Going from a 4\,MB table (in L3) to DRAM costs 2.3--5.8$\times$ for every query and
layout, more than any layout difference for wide queries.
\end{itemize}

\subsection{PAX page size}
\label{sec:res:pax}

\begin{figure}[t]
\centering
\fig{pax_pagesize}
\caption{PAX latency vs page size: plain packing, 64 bytes of padding after each minipage,
and an odd-line minipage stride (\code{--pax-pad auto}). Horizontal lines: DSM with the same
query and strategy. The predicate on $c_0$ passes 1\% or all rows.}
\label{fig:pax}
\end{figure}

\Cref{fig:pax} varies the PAX page size from 1\,KB to 4\,MB. A minipage is $P/16$ bytes: 64
bytes at 1\,KB, 1\,KB at the default 16\,KB, 256\,KB at 4\,MB.
\begin{itemize}
\item \textbf{One attribute} (SUM(c0)) gets faster as pages grow: 20.1\ms{} at 1\,KB, 4.9 at
16\,KB, 2.8 at 64\,KB and 1.2 at 4\,MB, which equals DSM. At 16\,KB the excess over DSM is
3.7\ms{} over 31{,}250 pages, about 120\,ns per page, roughly one DRAM latency.
\item \textbf{\code{SELECT *} at 1\%.} Late falls from 12.9\ms{} at 16\,KB to 6.5\ms{} at
256\,KB, reaching DSM late (7.1\ms). Early with plain packing stays at 38--66\ms{} at every
size, 1.7--2.9$\times$ DSM early (22.7\ms), and gets \emph{worse} from 64\,KB (65.6\ms) as
pages grow; with an odd-line minipage stride it drops to 23--28\ms.
\item \textbf{\code{SELECT *} of every row.} Late is fastest with tiny pages (71--72\ms{} at
1--2\,KB, 86--101\ms{} from 16\,KB): a page of 8 tuples keeps a vector's output tile in L1.
\item \textbf{64 bytes of padding} per minipage changes nothing from 16\,KB up and wastes
space below.
\end{itemize}
So the PAX page size trades the prefetcher, which wants long minipages, against
reconstruction locality, which wants small pages. The 8--16\,KB pages of disk-based PAX are on
the wrong side of it for scans in memory. \Cref{sec:analysis:cache} explains both effects and
\cref{sec:tuning} fixes them.
